% Figure F: Gradient magnitude decay — bar chart
\begin{tikzpicture}
\begin{axis}[
  width=10cm, height=5.5cm,
  ybar,
  bar width=0.5cm,
  xlabel={Timestep (distance from loss)},
  ylabel={Gradient magnitude},
  xtick={1,2,3,4,5,6,7,8},
  xticklabels={$t{=}T$,$t{=}T{-}1$,$t{=}T{-}2$,$t{=}T{-}3$,
               $t{=}T{-}4$,$t{=}T{-}5$,$t{=}T{-}6$,$t{=}T{-}7$},
  xticklabel style={font=\scriptsize, rotate=30, anchor=east},
  yticklabel style={font=\scriptsize},
  xlabel style={font=\small},
  ylabel style={font=\small},
  ymajorgrids=true,
  grid style={dashed, thwsgrey!60},
  axis line style={thwsgrey!80},
  ymin=0, ymax=1.15,
  enlarge x limits=0.08,
  /pgf/number format/1000 sep={},
]
\addplot[fill=thwsblue!60, draw=thwsblue] coordinates {
  (1, 1.0)
  (2, 0.65)
  (3, 0.42)
  (4, 0.27)
  (5, 0.18)
  (6, 0.11)
  (7, 0.07)
  (8, 0.05)
};
\end{axis}

% Annotation
\node[font=\scriptsize, align=center, color=carnelian] at (7.5, 3.5)
  {\alert{early inputs receive}\\nearly zero gradient};
\draw[->, >=stealth, color=carnelian, thick] (7.5, 3.1) -- (6.2, 1.5);

\end{tikzpicture}
